\documentclass{article}
\usepackage[UTF8]{ctex}
\usepackage{amsmath,amssymb,amsfonts}
\usepackage{enumitem}
\usepackage{geometry}
\geometry{a4paper, margin=1in}
% === 错题本排版优化（不参与题目标识） ===
\setlength{\parindent}{0pt}
\setlength{\parskip}{0.35em plus 0.10em minus 0.08em}
\setlength{\emergencystretch}{3em}
\setlist{itemsep=0.25em, topsep=0.40em, parsep=0pt}
\allowdisplaybreaks[2]
% === 错题本排版优化结束 ===
\begin{document}
\section*{湖南大学2026年数学分析真题}
\section*{143．湖南大学2026年数学分析试卷}

\subsection{题目1}
1．试求极限 \(\lim _{n \rightarrow \infty} \sum_{i=1}^{n} \frac{\sin \frac{i}{n} \pi}{n+\frac{i}{n}}\) ．

\subsection{题目2}
2．设 \(f(x)\) 在 \([0,1]\) 上连续，在 \((0,1)\) 上二阶可导，\(f(0) f(1)>0\) 且 \(f^{\prime \prime}(x)>0, \int_{0}^{1} f(x) \mathrm{d} x=0\) ．

\begin{itemize}
  \item[（1）] 证明：\(f(x)\) 在 \((0,1)\) 上恰有两个零点．
  \item[（2）] 证明：存在 \(c \in(0,1)\) ，使得 \(f^{\prime}(c)=\int_{0}^{c} f(x) \mathrm{d} x\) ．
\end{itemize}

\subsection{题目3}
3．解答如下问题：

\begin{itemize}
  \item[（1）] 证明：\(f(x, y)=\sqrt{x^{2}+y^{2}}\) 在 \(\mathbb{R}^{2}\) 上一致连续。
  \item[（2）] 证明：\(h(x, y)=\left(x^{2}+y^{2}\right)^{\frac{3}{2}}\) 在 \(\mathbb{R}^{2}\) 上非一致连续．
\end{itemize}

\subsection{题目4}
4．设 \(\left(x_{0}, y_{0}\right) \in \mathbb{R}^{2}, f(x, y)\) 在（ \(x_{0}, y_{0}\) ）处连续，\(g(x, y)\) 在（ \(x_{0}, y_{0}\) ）处可微且 \(g\left(x_{0}, y_{0}\right)=0\) ，证明： \(f(x, y) g(x, y)\) 在 \(\left(x_{0}, y_{0}\right)\) 处可微．

\subsection{题目5}
5．设 \(f \in C^{2}[0,1], f(0)=f(1)=0, f(x)\) 在（0 ，1）上无零点，证明： \(\int_{0}^{1}\left|\frac{f^{\prime \prime}(x)}{f(x)}\right| \mathrm{d} x>4\) ．

\subsection{题目6}
6．设 \(k>1\) ，且 \(k y-\sin y+x \cos y=0\) ．

\begin{itemize}
  \item[（1）] 证明：该方程在 \(|x|<k-1,|y|<+\infty\) 上确定唯一一个连续可导解 \(y=y(x)\) 。
  \item[（2）] 求 \(\lim _{x \rightarrow 0} \frac{y(x)}{x}\) ．
\end{itemize}

\subsection{题目7}
7．解答如下问题：

\begin{itemize}
  \item[（1）] 证明：\(\sum_{n=0}^{\infty} \int_{0}^{x} t^{n} \sin \pi t \mathrm{~d} t=\int_{0}^{x} \frac{\sin \pi t}{1-t} \mathrm{~d} t\) ，其中 \(0 \leq x \leq 1\) 。
  \item[（2）] 证明：\(\sum_{n=0}^{\infty} \int_{0}^{x} t^{n} \sin \pi t \mathrm{~d} t\) 在［0，1］上一致收敛。
\end{itemize}

\subsection{题目8}
8．设有收敛级数 \(\sum_{n=1}^{\infty} a_{n}=S\) 。

\begin{itemize}
  \item[（1）] 求 \(\lim _{n \rightarrow \infty} \frac{a_{1}+2 a_{2}+\cdots+n a_{n}}{n}\) ．
  \item[（2）] 证明：\(\sum_{n=1}^{\infty} \frac{a_{1}+2 a_{2}+\cdots+n a_{n}}{n(n+1)}\) 收敛，并求其和。
\end{itemize}

\subsection{题目9}
9．求 \(I=\iiint_{D}(x+y) \mathrm{d} x \mathrm{~d} y \mathrm{~d} z\) ，其中 \(D\) 为 \(x=0, x=1, x^{2}+1=\frac{y^{2}}{a^{2}}+\frac{z^{2}}{b^{2}}(a, b>0)\) 围成的区域．
\end{document}
